\section{Limitations}
\label{sec:limitations}
% Jujur & spesifik — reviewer menghargai ini. Sumber: caveat di findings.

\begin{itemize}
  \item \textbf{The map replicates across families; the causal results
        do not yet.} The representational findings are replicated on
        three checkpoints of a second model family
        (Sec.~\ref{sec:crossmodel}), but the patching experiments---and
        with them the map--use dissociation---have only been run on
        Mistral-7B. The output-side near-invariance does replicate on an
        instruction-tuned Mistral under its chat template
        (Sec.~\ref{sec:ceiling}), so it is not a base-model artifact,
        but it remains a claim about the letter-probability interface.
  \item \textbf{Letter-probability read-out.} Output distributions are
        read from option-letter logits after \texttt{Answer:}---the
        standard silicon-sampling interface, but known to diverge from
        free-text answers \citep{myanswerisc2024}. Our output-side claims
        (Sec.~\ref{sec:gap}) therefore scope to this interface: they say
        that identity barely reaches the letter-probability read-out, not
        that it fails to reach free-text generation, which we did not
        measure. The representation-side claims (Sec.~\ref{sec:act2}) do
        not depend on the read-out interface at all.
  \item \textbf{Single-institution, US-centric ground truth.} All
        distances are measured against Pew ATP; cross-institution
        transfer (e.g., scoring the same map against GSS-derived
        inter-group distances, to rule out one survey house's question
        style) and cross-country replication (e.g., WVS, with a known
        language confound) are future work.
  \item \textbf{Correlational map, partial causal probe.} RSA fidelity is
        correlational; the patching program (Sec.~\ref{sec:v2},
        Sec.~\ref{sec:sweep}) addresses use, not the origin of the map.
        Our dissociation claim rests on six attribute types---enough for
        two unambiguous counterexamples, not enough to estimate the
        fidelity--causality relationship with precision.
  \item \textbf{\starhead{} discovery provenance.} The head was first
        noticed in per-type top-10 lists before being fixed for the
        all-type test; strongest confirmation requires a second model /
        preregistered replication.
  \item \textbf{Sparse cells.} Some intersectional cells have thin survey
        samples; NaN pairs were excluded (Appendix~\ref{app:data}).
  \item \textbf{Twelve pairs per type bound the causal power.} With
        cluster-robust inference the causal experiments have 12 effective
        units per type: enough for one clear result and for direct
        between-type contrasts, not enough to certify a sweep winner
        against 31 rivals or to estimate the fidelity--causality
        relationship precisely ($n{=}6$ types). Scaling pairs, not
        questions, is the binding constraint for follow-up work.
  \item \textbf{The lexical control is one encoder.} The partial-fidelity
        control (Sec.~\ref{sec:lexical}) uses a single small sentence
        encoder as the lexical baseline; static word vectors cannot
        represent several of our value strings (e.g.\ age ranges). A
        stronger lexical model would only make the control stricter.
\end{itemize}

\section*{Ethics Statement}
Simulating the opinions of demographic groups is dual-use. It supports
legitimate survey-methodology research---pilot design, non-response
modelling, coverage checks---but the same capability can be used to
manufacture apparently representative public opinion, or to reinforce
stereotypes by presenting a model's guess about a group as that group's
view.

Our findings cut against the permissive reading. We show that current
LLM outputs should \emph{not} be treated as substitutes for human
respondents: under natural prompting the model's answers move by less
than 2\% of their error when the entire demographic identity is replaced,
and the small movement that exists is not reliably toward the target
group's truth. The failure is also unequal. Race-based attributes are the
weakest and the most prompt-fragile everywhere we measured, so simulated
outputs for racial groups are the least trustworthy exactly where
misuse would be most harmful. We therefore caution explicitly against
policy-adjacent uses of simulated group opinions.

We also note a risk specific to interpretability results of this kind: a
faithful location is a natural target for steering, and our own
Sec.~\ref{sec:v2} shows that naive steering through such a location can
push predictions \emph{away} from the truth in low-fidelity types. Read
the map before trusting it as a lever.

All data are public, aggregated Pew ATP response distributions with no
personally identifiable information; no human subjects were recruited for
this work.
